\documentclass[11pt]{article}
% styles.tex
\usepackage[utf8]{inputenc}
\usepackage[T1]{fontenc}
\usepackage[spanish, es-theorem]{babel}
\usepackage{amsmath, amssymb, amsthm}
\usepackage{mathrsfs}
\usepackage{geometry}
\geometry{a4paper, margin=2.8cm, headheight=14pt}
\usepackage{hyperref}
\hypersetup{colorlinks=true, linkcolor=black, citecolor=black, urlcolor=black}
\usepackage{fancyhdr}
\usepackage{titlesec}
\usepackage{cite}
\usepackage{newtxtext,newtxmath}
\usepackage{graphicx} 
\usepackage{float}

% Entornos de teoremas estilo Apostol
\newtheorem{theorem}{Teorema}[section]
\newtheorem{lemma}[theorem]{Lema}
\newtheorem{remark}[theorem]{Observación}
\newtheorem{definition}[theorem]{Definición}
\newtheorem{proposition}[theorem]{Proposición}
\newtheorem{corollary}[theorem]{Corolario}
\newtheorem{axiom}[theorem]{Axioma}
\newcommand{\norm}[1]{\left\|#1\right\|}

\usepackage{listings}
\usepackage{xcolor}

% Definición de colores para el código
\definecolor{codegreen}{rgb}{0,0.6,0}
\definecolor{codegray}{rgb}{0.5,0.5,0.5}
\definecolor{codepurple}{rgb}{0.58,0,0.82}
\definecolor{backcolour}{rgb}{0.95,0.95,0.92}

\lstdefinestyle{pythonstyle}{
    backgroundcolor=\color{backcolour},   
    commentstyle=\color{codegreen},
    keywordstyle=\color{magenta},
    numberstyle=\tiny\color{codegray},
    stringstyle=\color{codepurple},
    basicstyle=\ttfamily\scriptsize,
    breakatwhitespace=false,         
    breaklines=true,                 
    captionpos=b,                    
    keepspaces=true,                 
    numbers=left,                    
    numbersep=5pt,                  
    showspaces=false,                
    showstringspaces=false,
    showtabs=false,                  
    tabsize=4
}
\lstset{style=pythonstyle}

% Configuración de página
\pagestyle{fancy}
\fancyhf{}
\lhead{\small \textit{Operador de Integridad de Bisagra}}
\rhead{\small J. Knuttzen}
\cfoot{\thepage}

% Formato de títulos
\titleformat{\section}{\large\bfseries}{\thesection.}{0.5em}{}
\titleformat{\subsection}{\normalsize\bfseries}{\thesubsection.}{0.5em}{}


% --- TÍTULO ---
\title{\textbf{Formalismo de Integridad de Bisagra: \\ Una Aproximación Espectral y Topológica a la Estabilidad Geométrica y la Localización de Singularidades}}
\author{Joaquín Knuttzen \\ \small \textit{Investigador Independiente, Santa Fe, Argentina}}
\date{\today}

\begin{document}
\maketitle

\begin{abstract}
Se presenta el Operador de Integridad de Bisagra (IOB), un funcional escalar diseñado para cuantificar la preservación de la suavidad métrica en variedades diferenciables mediante la evaluación de residuos bilaterales simétricos. Este formalismo trasciende las limitaciones locales de la linealización de primer orden al mapear el estrés métrico hacia un dominio espectral (IOB-FFT). Para reducir la complejidad computacional inherente a la evaluación volumétrica de campos hiper-trascendentales, se introduce el Teorema del Flujo de Integridad, el cual proyecta este operador sobre integrales de contorno. Esta formulación perimetral fundamenta una arquitectura de partición espacial recursiva (IOB-QuadTree) que disminuye drásticamente la latencia de búsqueda al aislar singularidades evaluando exclusivamente las fronteras del espacio de fases. Adicionalmente, el operador se generaliza como un sensor dinámico local para sistemas continuos. La robustez del método se valida frente a patologías de fragmentación fractal y demuestra capacidad anticipatoria ante bifurcaciones inminentes en fluidodinámica, osciladores acoplados y transiciones de régimen en el atractor de Lorenz.
\end{abstract}

\vspace{0.5cm}
\noindent \textbf{Palabras clave:} Operador de Integridad de Bisagra, Estabilidad topológica, Exponentes de Lyapunov, Métodos sin Jacobiano, Transiciones críticas, Análisis de singularidades.

% -------- SECCIÓN 1: INTRODUCCIÓN --------
\section{Introducción}

El análisis de estabilidad geométrica y la localización de estados críticos han dependido históricamente de la viabilidad de linealizar campos vectoriales en la vecindad de una singularidad. En sistemas caracterizados por clústeres de alta densidad u horizontes de bifurcación, la aproximación mediante series de Taylor truncadas exhibe una degradación sistemática. Esta pérdida de fidelidad está dictada por la contracción del radio de convergencia analítico impuesto por las singularidades adyacentes.

Aunque métodos como la continuación homotópica \cite{sommese2005numerical} han resuelto con notable eficacia la extracción de raíces para sistemas estrictamente polinomiales, su aplicabilidad se ve severamente restringida frente a funciones trascendentales o campos con singularidades esenciales. En este dominio, los esquemas iterativos enfrentan la fragmentación de las cuencas de convergencia en estructuras fractales \cite{smale1981fundamental, kravanja2000computing}, lo que subordina la localización de soluciones a una búsqueda de fuerza bruta o a la extrema precisión de una semilla inicial.

Este trabajo propone una reformulación de la métrica de estabilidad mediante el \textbf{Operador de Integridad de Bisagra (IOB)}, el cual evalúa la topología del espacio de fases no como un gradiente puntual estrictamente local, sino como la coherencia de un tejido métrico sometido a deformación bilateral. 

\newpage
\subsection{Limitaciones de la Información de Primer Orden y la Alternativa Perimetral}

Sea un campo $f$ evaluado en una vecindad de $\mathbf{x}_0$. La aproximación de Fréchet presupone que la topología de la variedad local es asintóticamente plana congruente con su espacio tangente \cite{lee2012smooth}. Sin embargo, en dominios donde el tensor diferencial colapsa ($\|df_{\mathbf{x}_0}\| \to 0$) o la cizalladura métrica diverge, los términos residuales no lineales dominan la evolución del sistema. 

El operador IOB cuantifica esta divergencia normalizando el residuo cuadrático bilateral respecto a la escala espacial de muestreo. Esta operación transforma una divergencia algebraica asintótica en un observable escalar estrictamente acotado, denominado \textit{índice de integridad}. De mayor relevancia analítica, al proyectar esta medida sobre una curva cerrada $\gamma$, el operador permite deducir analíticamente la presencia de singularidades internas interrogando exclusivamente la frontera, eludiendo así la penalización computacional de la integración volumétrica.

A lo largo de este artículo, se formalizarán las propiedades fundamentales del IOB (Sección \ref{sec:fundamentos}) y se proyectará su aplicación en dos frentes analíticos mayores. Por un lado (Sección \ref{sec:flujo}), se demostrará su eficacia en la resolución simultánea de raíces en el plano complejo mediante el análisis global de Fourier (IOB-FFT) y se formulará un algoritmo de partición espacial (IOB-QuadTree) impulsado por el cálculo del flujo de integridad perimetral. Por otro lado (Sección \ref{sec:loc_universal}), se extenderá su formulación al dominio temporal como un sensor Euleriano-Lagrangiano continuo, operando como un predictor de crisis topológicas. Finalmente, en la Sección \ref{sec:dinamica_campos}, el formalismo se someterá a validación empírica en escenarios de estrés extremo: desde el aislamiento de resonancias hiper-trascendentales y el caos mecánico, hasta la detección de ondas de choque en fluidodinámica compresible.

% -------- SECCIÓN 2: FUNDAMENTOS IOB --------
\section{Fundamentos y Propiedades del Operador IOB}
\label{sec:fundamentos}

La caracterización de la estabilidad geométrica de una aplicación diferenciable $f: \Omega \subseteq \mathbb{R}^n \to \mathbb{R}^m$ requiere, tradicionalmente, la evaluación de su estructura diferencial de primer orden. No obstante, en la vecindad de singularidades topológicas, el diferencial de Fréchet $df_{\mathbf{x}}$ deja de ser un descriptor suficiente de la variedad local. Introducimos el \textbf{Operador de Integridad de Bisagra (IOB)} como un funcional escalar no lineal que cuantifica la desviación estricta de la variedad respecto a su hiperplano tangente, explotando las propiedades de simetría de un muestreo bilateral.

\subsection{Definición del Residuo Bilateral}

Sea $\mathbf{x} \in \Omega$ un punto de evaluación y $\mathbf{u} \in \mathbb{S}^{n-1}$ un vector unitario que define una dirección de perturbación. Dada una escala espacial isotrópica $L \in \mathbb{R}^+$, definimos los residuos bilaterales $\mathbf{R}_+$ y $\mathbf{R}_- \in \mathbb{R}^m$ como la desviación exacta de la aproximación lineal de Fréchet evaluada en una vecindad simétrica a lo largo del eje $\mathbf{u}$:
\begin{equation}
    \mathbf{R}_{\pm}(\mathbf{x}, \mathbf{u}, L) = f(\mathbf{x} \pm L\mathbf{u}) - [f(\mathbf{x}) \pm L \cdot df_{\mathbf{x}}(\mathbf{u})] \,.
\end{equation}
Bajo este formalismo, la coordenada pivot $\mathbf{x}$ actúa como una bisagra infinitesimal. Para una aplicación estrictamente afín, la norma euclidiana de ambos residuos espaciales es idénticamente nula en todo $\Omega$. Consecuentemente, la emergencia de residuos espaciales direccionales indica una ruptura de la planitud isométrica local impuesta por la escala $L$.

\begin{remark}[Aliasing Topológico y Escala]
Desde una perspectiva analítica, considerar a $L$ como un escalar estático global introduce el riesgo de \textit{aliasing topológico}. En dominios caracterizados por una alta densidad de singularidades (e.g., cuencas fractales), si la perturbación $L$ supera la semidistancia mínima entre raíces adyacentes, el operador promediará la curvatura de múltiples cuencas, fracasando en la resolución espectral. Para garantizar la inyectividad del operador, $L$ debe estar estrictamente acotado por $L < \frac{1}{2} \min_{i \neq j} \|r_i - r_j\|$, o bien subordinarse dinámicamente al gradiente local mediante un esquema adaptativo del tipo $L(\mathbf{x}) \propto \|\nabla \Phi(\mathbf{x})\|^{-1}$.
\end{remark}

\subsection{El Índice de Integridad Direccional $\mathcal{I}$}

Para cuantificar la preservación métrica local de manera independiente a la escala absoluta de la función, definimos el índice de integridad $\mathcal{I}$ como una métrica adimensional. Utilizando la norma euclidiana $\|\cdot\|_2$ en $\mathbb{R}^m$ y normalizando el estrés respecto a la magnitud de la deformación lineal esperada, el operador se define mediante el funcional:
\begin{equation}
    \label{eq:indice_integridad}
    \mathcal{I}(\mathbf{x}, \mathbf{u}, L) = 1 - \frac{\|\mathbf{R}_+\|_2^2 + \|\mathbf{R}_-\|_2^2}{2L^2 (\|df_{\mathbf{x}}(\mathbf{u})\|_2^2 + \epsilon)} \,,
\end{equation}
donde $\epsilon = \varepsilon_{\text{mach}} \max(1, \|df_{\mathbf{x}}(\mathbf{u})\|_2^2)$ actúa como un regularizador aritmético vinculado al épsilon de máquina ($\varepsilon_{\text{mach}}$). Este término previene divergencias espurias en puntos estacionarios nulos sin alterar el límite asintótico del funcional. El factor fraccionario representa la densidad relativa de energía de deformación no lineal proyectada sobre el eje $\mathbf{u}$. Por construcción, $\mathcal{I} \approx 1$ si y solo si la variedad local coincide con su espacio tangente direccional.

\subsection{Comportamiento Asintótico y Simetría de Cancelación}

La validación del IOB como descriptor diferencial riguroso exige establecer su convergencia asintótica y su vinculación directa con los tensores métricos de orden superior. La arquitectura bilateral del muestreo otorga al operador una precisión asintótica superior a la estimación asimétrica estándar (como las diferencias finitas \textit{forward}).

\begin{lemma}[Convergencia de Fréchet]
Sea $f$ diferenciable en el sentido de Fréchet en $\mathbf{x} \in \Omega$. Para toda dirección $\mathbf{u} \in \mathbb{S}^{n-1}$, el índice de integridad satisface el límite asintótico plano:
\begin{equation}
    \lim_{L \to 0} \mathcal{I}(\mathbf{x}, \mathbf{u}, L) = 1 \,.
\end{equation}
\end{lemma}

\begin{proof}
Por definición analítica de la diferenciabilidad de Fréchet en variedades euclidianas \cite{lee2012smooth}, la norma de los residuos satisface $\|\mathbf{R}_{\pm}\|_2 = o(L)$ cuando $L \to 0$. Sustituyendo esta condición en la Ecuación \eqref{eq:indice_integridad}, el término de estrés métrico decae proporcionalmente a $(o(L)/L)^2 \to 0$, garantizando la recuperación estricta de la isometría afín en el límite infinitesimal.
\end{proof}

\begin{theorem}[Aislamiento de Curvatura Intrínseca y Decaimiento Par] \label{thm:aislamiento_curvatura}
Si $f \in C^3$ en una vecindad convexa de $\mathbf{x}$, el estrés métrico exhibe un decaimiento estrictamente par respecto a la escala $L$. Este fenómeno aísla el diferencial de segundo orden suprimiendo analíticamente la perturbación asimétrica de tercer orden (cizalladura cúbica), operando bajo el mismo principio de cancelación de orden impar que fundamenta los esquemas de diferencias finitas centrales \cite{burden2015numerical}:
\begin{equation}
    \mathcal{I}(\mathbf{x}, \mathbf{u}, L) = 1 - \frac{L^2}{4} \|d^2 f_{\mathbf{x}}(\mathbf{u}, \mathbf{u})\|_2^2 + \mathcal{O}(L^4) \,.
\end{equation}
\end{theorem}

\begin{proof}
Considerando la expansión de Taylor en torno a la coordenada pivot $\mathbf{x}$, los residuos bilaterales toman la forma:
\begin{equation*}
    \mathbf{R}_{\pm} = \frac{L^2}{2} d^2 f_{\mathbf{x}}(\mathbf{u}, \mathbf{u}) \pm \frac{L^3}{6} d^3 f_{\mathbf{x}}(\mathbf{u}, \mathbf{u}, \mathbf{u}) + \mathcal{O}(L^4) \,.
\end{equation*}
Al evaluar la suma de las normas cuadráticas en el numerador del IOB, $\|\mathbf{R}_+\|_2^2 + \|\mathbf{R}_-\|_2^2$, el producto interno cruzado entre los términos de deformación simétrica ($\mathcal{O}(L^2)$) y asimétrica ($\mathcal{O}(L^3)$) presenta signos estrictamente opuestos. Esta paridad induce una cancelación idéntica del término cúbico:
\begin{equation*}
    2 \left\| \frac{L^2}{2} d^2 f_{\mathbf{x}} \right\|_2^2 + \mathcal{O}(L^6) = \frac{L^4}{2} \|d^2 f_{\mathbf{x}}\|_2^2 + \mathcal{O}(L^6) \,.
\end{equation*}
Al dividir este sumatorio por el factor de normalización $2L^2$ inherente a la Ecuación \eqref{eq:indice_integridad}, se obtiene la expresión del teorema. Esta cancelación geométrica del tensor de tercer orden es el mecanismo fundamental que confiere al IOB una relación señal-ruido excepcionalmente alta y un error de truncamiento de orden $\mathcal{O}(L^4)$.
\end{proof}

\subsection{Generalización: El Operador de Integridad Isótropo}

El índice direccional $\mathcal{I}$ es insuficiente para describir la variedad globalmente, dado que es direccionalmente dependiente. Para extraer un invariante escalar intrínseco del tensor métrico en el punto $\mathbf{x}$, definimos el \textbf{Operador de Integridad Isótropo} $\overline{\mathcal{I}}$, integrando la densidad de estrés direccional sobre la medida de Lebesgue normalizada de la hiperesfera $\mathbb{S}^{n-1}$:
\begin{equation}
    \overline{\mathcal{I}}(\mathbf{x}, L) = 1 - \frac{1}{|\mathbb{S}^{n-1}|} \int_{\mathbb{S}^{n-1}} Q(\mathbf{x}, \mathbf{u}, L) \, d\mathbf{u} \,,
\end{equation}
donde $Q(\mathbf{x}, \mathbf{u}, L)$ representa la densidad direccional de energía de deformación elástica (el término fraccionario sustractivo de la Ecuación \eqref{eq:indice_integridad}).

Para viabilizar su implementación algorítmica, la integral continua sobre la variedad $\mathbb{S}^{n-1}$ se proyecta mediante una cuadratura esférica simétrica. Específicamente, para el caso bidimensional de localización en el plano complejo ($n=2$), introducimos un conjunto de $M$ vectores ortonormales uniformemente distribuidos $\mathbf{u}_k = (\cos \theta_k, \sin \theta_k)$ con $\theta_k = 2\pi k / M$. El operador isótropo computacional discreto se formula como:
\begin{equation}
    \overline{\mathcal{I}}_{comp}(\mathbf{x}, L) = 1 - \frac{1}{M} \sum_{k=1}^{M} \frac{\|\mathbf{R}_+(\mathbf{x}, \mathbf{u}_k)\|_2^2 + \|\mathbf{R}_-(\mathbf{x}, \mathbf{u}_k)\|_2^2}{2L^2 \left(\|df_{\mathbf{x}}(\mathbf{u}_k)\|_2^2 + \epsilon\right)} \,.
\end{equation}
Para preservar la invariancia de cancelación detallada en el Teorema \ref{thm:aislamiento_curvatura}, es condición necesaria y suficiente que el número de nodos sea $M \geq 2$ pares diametralmente opuestos. En mallas de discretización estándar, un muestreo en cruz ortogonal ($M=4$) captura la norma de Frobenius del Hessiano matricial, maximizando la eficiencia computacional con complejidad $\mathcal{O}(1)$ por coordenada espacial.


% -------- SECCIÓN 3: INTEGRACIÓN TOPOLÓGICA --------
\section{Integración Topológica y el Teorema del Flujo de Integridad}
\label{sec:flujo}

El análisis complejo clásico establece \cite{ahlfors1979complex}, mediante el Principio del Argumento y el Teorema de los Residuos de Cauchy, que la evaluación de una integral de contorno sobre la derivada logarítmica de una función analítica cuantifica estrictamente el número de singularidades (raíces y polos) encerradas en su interior. No obstante, la dependencia de la derivada analítica $f'(z)$ convierte a este enfoque en una herramienta computacionalmente frágil y sujeta a inestabilidades numéricas severas frente a funciones hiper-trascendentales o dinámicas de ``caja negra''.

En esta sección, se reformula el conteo topológico proyectando el Operador de Integridad de Bisagra hacia integrales de contorno. Esta maniobra prescinde por completo del cálculo derivativo explícito y de la evaluación de la fase compleja, transformando el problema en un balance de flujos escalares de deformación métrica.

\subsection{El Potencial de Fase y el Flujo Gradiente}

Sea $f: \Omega \subset \mathbb{C} \to \mathbb{C}$ una función analítica en un dominio cerrado $D \subset \Omega$, cuya frontera es una curva cerrada simple $\gamma$. Definimos el potencial escalar de fase como $\Phi(z) = \ln|f(z)|$.

En virtud de las ecuaciones de Cauchy-Riemann, la derivada direccional normal del logaritmo de la magnitud ($\nabla \Phi \cdot \mathbf{\hat{n}}$) es estrictamente equivalente a la derivada tangencial del argumento complejo ($\frac{\partial}{\partial s} \arg(f)$). Al integrar esta equivalencia direccional sobre el contorno simple $\gamma$, el Teorema de la Divergencia de Gauss en el plano \cite{marsden2003vector} permite aislar la carga topológica neta $\mathcal{N}$ (la diferencia exacta entre el número de ceros $Z$ y polos $P$, contando su multiplicidad):
\begin{equation}
    \label{eq:teorema_divergencia}
    \mathcal{N} = Z - P = \frac{1}{2\pi} \oint_{\gamma} \nabla \Phi \cdot \mathbf{\hat{n}} \, ds \,.
\end{equation}
Aunque matemáticamente elegante, la evaluación algorítmica continua de $\nabla \Phi$ arrastra las mismas penalizaciones numéricas que la derivada analítica iterativa, limitando su utilidad práctica en mallas discretizadas.

\subsection{El Teorema del Flujo de Integridad (Formulación de Bisagra)}

Para evadir el cálculo diferencial local, aplicamos el principio de perturbación bilateral del IOB (desarrollado en la Sección \ref{sec:fundamentos}) sobre la curva $\gamma$. Evaluamos el potencial $\Phi$ mediante un muestreo ortogonal a la frontera, separado por una distancia infinitesimal isotrópica $L$, orientada en la dirección del vector normal saliente $\mathbf{\hat{n}}$.

\begin{theorem}[Flujo de Integridad de Bisagra] \label{thm:teorema_flujo}
Sea $L \in \mathbb{R}^+$ el radio de muestreo bilateral. La aproximación discreta de simetría par para la carga topológica encerrada por la curva $\gamma$ está dada por el operador de contorno $\mathcal{N}_{IOB}$:
\begin{equation}
    \label{eq:flujo_bisagra}
    \mathcal{N}_{IOB} = \frac{1}{4\pi L} \oint_{\gamma} \Big( \ln|f(z + L\mathbf{\hat{n}})| - \ln|f(z - L\mathbf{\hat{n}})| \Big) ds \,.
\end{equation}
Este operador asintótico converge al número entero de la carga topológica $\mathcal{N}$ suprimiendo idénticamente las perturbaciones asimétricas espurias de la variedad local.
\end{theorem}

\begin{proof}
Aproximamos la derivada direccional normal $\nabla \Phi \cdot \mathbf{\hat{n}}$ presente en la Ecuación \eqref{eq:teorema_divergencia} mediante un esquema de discretización de diferencia central simétrica \cite{burden2015numerical}:
\begin{equation*}
    \nabla \Phi \cdot \mathbf{\hat{n}} = \frac{\Phi(z + L\mathbf{\hat{n}}) - \Phi(z - L\mathbf{\hat{n}})}{2L} + \mathcal{O}(L^2) \,.
\end{equation*}
Sustituyendo $\Phi(z) = \ln|f(z)|$ e insertando esta aproximación en la integral de la Ecuación \eqref{eq:teorema_divergencia}, el factor divisor $2L$ se factoriza y multiplica por el coeficiente $2\pi$ inherente al Teorema de Gauss, obteniendo el escalar normalizador $4\pi L$ de la Ecuación \eqref{eq:flujo_bisagra}. 

Análogamente a la demostración de invariancia del Teorema \ref{thm:aislamiento_curvatura} (Sección \ref{sec:fundamentos}), la simetría par del muestreo bilateral aniquila de manera idéntica los términos de error de orden impar de la expansión de Taylor. Por lo tanto, el error de truncamiento del integrando decae como $\mathcal{O}(L^2)$. Esta cancelación analítica permite operar con radios $L$ macroscópicos sin comprometer la convergencia hacia el cuásar numérico entero que representa la carga topológica.
\end{proof}


% -------- SECCIÓN 4: LOCALIZACIÓN Y EXTRACCIÓN --------
\newpage
\section{Localización Universal de Singularidades en el Plano Complejo}
\label{sec:loc_universal}

La extracción determinista de raíces para funciones analíticas no polinomiales (trascendentales o con singularidades esenciales) constituye un desafío abierto en el análisis numérico. A diferencia de los sistemas puramente algebraicos, donde métodos como la continuación homotópica resultan eficientes, el paisaje topológico de funciones trascendentales fragmenta las cuencas de convergencia de los esquemas iterativos estándar, subordinándolos al caos determinista y a la ceguera espacial.

En esta sección, se proyecta el Teorema del Flujo de Integridad \ref{thm:teorema_flujo} (Sección \ref{sec:flujo}) hacia el plano complejo $\mathbb{C}$. Se proponen dos arquitecturas algorítmicas complementarias: un esquema de partición recursiva de alta eficiencia computacional (IOB-QuadTree) para la localización aislada, y un mapeo espectral continuo (IOB-FFT) para el análisis de clústeres densos.

\subsection{Algoritmo de Exclusión Topológica (IOB-QuadTree)}

Los algoritmos de búsqueda global tradicionales evalúan funciones sobre mallas bidimensionales exhaustivas, incurriendo en un costo computacional de $\mathcal{O}(M^2)$. Para quebrar esta asimetría, introducimos una estructura de datos de partición espacial recursiva impulsada por el operador de contorno $\mathcal{N}_{IOB}$.

Sea un dominio inicial acotado $D_0 \subset \mathbb{C}$ con frontera $\partial D_0$. El algoritmo IOB-QuadTree procede mediante la siguiente lógica de exclusión:
\begin{enumerate}
    \item \textbf{Evaluación Perimetral:} Se calcula numéricamente la carga topológica $\mathcal{N}_{IOB}$ (Ecuación \eqref{eq:flujo_bisagra}) evaluando el estrés métrico exclusivamente sobre el perímetro $\partial D_0$.
    \item \textbf{Criterio de Exclusión:} Si $\mathcal{N}_{IOB} < 1$, el Teorema del Flujo \ref{thm:teorema_flujo} garantiza la ausencia de singularidades en el interior de $D_0$. El dominio entero se descarta en tiempo $\mathcal{O}(1)$, eludiendo la evaluación de su volumen.
    \item \textbf{Subdivisión Recursiva:} Si $\mathcal{N}_{IOB} \geq 1$, el dominio se divide en cuatro sub-cuadrantes disjuntos $D_{k}^{(1)}$. El proceso se itera recursivamente hasta alcanzar una profundidad máxima predefinida, aislando cada singularidad en una celda espacial estrictamente convexa.
\end{enumerate}
Dado que el operador interroga las fronteras en lugar de las áreas, la complejidad asintótica de la búsqueda desciende a $\mathcal{O}(\kappa \cdot M)$, donde $\kappa$ es una constante vinculada a la densidad local de las raíces. 

\subsection{Límite de Saturación y Mapeo Espectral Global (IOB-FFT)}

En regímenes topológicos donde la semidistancia entre raíces adyacentes es inferior a la resolución límite del árbol de partición (clústeres de altísima densidad), el algoritmo IOB-QuadTree alcanza un estado de saturación. En tales escenarios, recurrimos a una evaluación espectral global del campo.

Para un dominio computacional discreto $\mathcal{D}_{comp}$, introducimos el potencial de estrés espectral $\tilde{\Phi}(z) = -\ln(|f(z)| + \epsilon)$, el cual invierte y regulariza el campo escalar original para que las raíces emerjan como singularidades positivas. Este potencial se procesa extrayendo su curvatura mediante el operador Laplaciano. Para maximizar la eficiencia asintótica a $\mathcal{O}(N \log N)$ sobre mallas densas, esta operación se ejecuta en el dominio de la frecuencia espacial utilizando la Transformada Rápida de Fourier ($\mathcal{F}$), definiendo el \textbf{Funcional de Estrés Espectral}:
\begin{equation}
    \Psi(z) = \left| \mathcal{F}^{-1} \left\{ \mathcal{F}\{\tilde{\Phi}(z) \cdot W(z)\} \cdot \|\mathbf{k}\|^2 \right\} \right| \,,
\end{equation}
donde $\|\mathbf{k}\|^2 = k_x^2 + k_y^2$ es la norma al cuadrado del vector de onda bidimensional y $W(z)$ es una ventana de apodización. Los máximos locales del funcional $\Psi(z)$ emergen como deltas de Dirac regularizadas, revelando la estructura interna del clúster topológico.

\subsection{Convergencia Confinada y Refinamiento Newtoniano}

Los métodos iterativos tradicionales, como Newton-Raphson, colapsan en las vecindades de las fronteras fractales de las cuencas de atracción, lo que habitualmente exige la adopción de esquemas amortiguados (\textit{Damped Newton}). 

Sin embargo, el IOB-QuadTree (y análogamente el IOB-FFT) garantiza de manera determinista que la coordenada semilla $z_{seed}^{(k)}$ devuelta por el algoritmo reside profundamente dentro de la vecindad estrictamente monotónica de la raíz. Consecuentemente, el fenómeno de la eyección fractal queda neutralizado desde la inicialización. El refinamiento final de la coordenada se ejecuta mediante la actualización canónica $\Delta z_{n} = -f(z_n)/f'(z_n)$, garantizando una convergencia cuadrática asintótica y eximiendo al integrador de estimar heurísticamente radios de confianza (\textit{trust regions}).

\subsection{Estabilidad Numérica y Mapeo Asintótico}

La implementación del formalismo IOB enfrenta un desafío aritmético crítico cuando la función objetivo $f(z)$ exhibe un crecimiento hiper-trascendental. Magnitudes que exceden el estándar IEEE 754 de doble precisión ($|f(z)| > 10^{308}$) inducen un desbordamiento catastrófico (overflow) previo a la evaluación logarítmica.

Para mitigar esta patología, postulamos un \textbf{Mapeo Atómico}. Para funciones que admiten una representación de la forma $f(z) = \exp(g(z)) - 1$, el potencial de estrés espectral $\tilde{\Phi}$ se evalúa directamente desde la estructura interna del exponente analítico $g(z)$:
\begin{equation}
    \tilde{\Phi}_{asym}(z) = 
    \begin{cases} 
    -\operatorname{Re}(g(z)) & \text{si } \operatorname{Re}(g(z)) \gg \eta \\
    -\ln | \text{expm1}(g(z)) | & \text{en otro caso}
    \end{cases}
\end{equation}
donde $\text{expm1}(\cdot)$ es el operador de precisión subyacente y $\eta$ es el umbral de desbordamiento de la arquitectura de hardware. Este isomorfismo proyecta el infinito analítico hacia un observable acotado. 

Incluso con un potencial estabilizado, las funciones con gradientes marcadamente explosivos inducen regiones de gradiente evanescente en el dominio transformado, donde el ruido de cuantización digital enmascara la ubicación de las raíces. Para preservar la resolución topológica de los picos de densidad, introducimos una \textbf{Saturación de Curvatura} mediante la transformación monotónica $\tilde{\Phi}_{norm}(z) = \tanh(\tilde{\Phi}_{asym}(z) / \alpha)$, donde $\alpha$ se deriva de la distribución intercuartílica de la métrica espacial. Dado que la transformación conserva la monotonicidad estricta, el operador preserva asintóticamente el lugar geométrico de las deltas de Dirac sin alterar las coordenadas de la singularidad.


% -------- SECCIÓN 5: SISTEMAS DINÁMICOS --------
\newpage
\section{Dinámica de Campos y Alertas Tempranas de Crisis Topológicas}
\label{sec:dinamica_campos}

El análisis de transiciones críticas en sistemas continuos autónomos $\dot{\mathbf{x}} = \mathbf{F}(\mathbf{x})$ se ha fundamentado históricamente en la teoría ergódica, utilizando descriptores como los Exponentes de Lyapunov en Tiempo Finito (FTLE) \cite{haller2015lagrangian, guckenheimer2013nonlinear}. No obstante, la evaluación de la divergencia exponencial exige la integración de las ecuaciones variacionales a lo largo de un horizonte temporal futuro $T$. Esta dependencia integral introduce una latencia ineludible: los FTLE confirman la inestabilidad retrospectivamente, cuando el sistema ya ha iniciado su divergencia.

Para la instrumentación predictiva en tiempo real (\textit{Early Warning Signals}) \cite{scheffer2009early}, es imperativo transicionar de una métrica de divergencia temporal a una métrica de deformación espacial estricta. En esta sección, proyectamos el Operador de Integridad Isótropo ($\overline{\mathcal{I}}$) como un tensor de monitoreo continuo para identificar el cizallamiento de la métrica del espacio de fases previo al colapso de una órbita hacia un atractor extraño o una extinción poblacional.

\subsection{El Sensor Dinámico Euleriano-Lagrangiano}

Sea una trayectoria $\mathbf{x}(t)$ evolucionando bajo el flujo del campo vectorial $\mathbf{F}: \Omega \subseteq \mathbb{R}^n \to \mathbb{R}^n$. Postulamos un esquema de medición donde la trayectoria física actúa como un observador Lagrangiano que, en cada instante $t$, evalúa el estrés Euleriano del campo vectorial circundante mediante el muestreo de bisagra.

\begin{definition}[Densidad de Estrés Métrico del Flujo]
Se define el estrés métrico dinámico $Q_d(t, L)$ como la evaluación escalar de la desviación isométrica del campo de velocidades en la vecindad de la coordenada instantánea $\mathbf{x}(t)$:
\begin{equation}
    \label{eq:estres_flujo}
    Q_d(t, L) = \frac{1}{|\mathbb{S}^{n-1}|} \int_{\mathbb{S}^{n-1}} \frac{\|\mathbf{R}_+^{\mathbf{F}}\|_2^2 + \|\mathbf{R}_-^{\mathbf{F}}\|_2^2}{2L^2 (\|d\mathbf{F}_{\mathbf{x}(t)}(\mathbf{u})\|_2^2 + \epsilon)} \, d\mathbf{u} \,,
\end{equation}
donde los residuos de bisagra del campo dinámico se expresan como:
\begin{equation*}
    \mathbf{R}_{\pm}^{\mathbf{F}} = \mathbf{F}(\mathbf{x}(t) \pm L\mathbf{u}) - [\mathbf{F}(\mathbf{x}(t)) \pm L \cdot d\mathbf{F}_{\mathbf{x}(t)}(\mathbf{u})] \,.
\end{equation*}
\end{definition}

Físicamente, el funcional $Q_d(t, L)$ cuantifica la curvatura geométrica estricta de las líneas de corriente. Cuando una trayectoria se aproxima a una separatriz u horizonte de bifurcación, la no-linealidad espacial se dispara, forzando a $Q_d(t)$ a exhibir un crecimiento divergente dictado por la contracción topológica del embudo de fase.

\subsection{Detección Causal y el Funcional Predictivo}
\label{sec:funcional_pred}
A diferencia de los regímenes analíticos estacionarios tratados en la Sección \ref{sec:loc_universal}, la predicción de fallas en sistemas físicos exige una estricta observancia del principio de causalidad computacional. La detección de la crisis topológica no puede depender de transformadas globales, sino de la historia inmediata del estrés.

Definimos el \textbf{Funcional Predictivo de Crisis} $\Psi_c(t)$ como la aceleración instantánea de la deformación métrica. Para garantizar la causalidad y suprimir el ruido micro-estructural inherente al integrador numérico, se evalúa la segunda derivada temporal sobre un perfil logarítmico suavizado de $Q_d(t)$:
\begin{equation}
    \Psi_c(t) = \left| \frac{d^2}{dt^2} \left( \ln(Q_d(t) + \epsilon) * w(t) \right) \right| \,,
\end{equation}
donde $w(t)$ es un núcleo de convolución causal continuo (ventana deslizante asimétrica con soporte en $[t-\tau, t]$). El parámetro $\tau$ representa el horizonte de memoria del filtro, el cual se sintoniza empíricamente en función de la autocorrelación intrínseca de la serie temporal para suprimir fluctuaciones transitorias sin enmascarar la dinámica subyacente. El operador de segunda derivada actúa, así, como un detector hiper-sensible de los puntos de inflexión en la compresión del tejido métrico.

\begin{theorem}[Pre-Alerta de Inestabilidad Dinámica]
\label{thm:pre_alerta}
Sea un sistema no lineal $\dot{\mathbf{x}} = \mathbf{F}(\mathbf{x})$ cuya trayectoria $\mathbf{x}(t)$ transita hacia una bifurcación local en el instante $t_{crit}$. Asumiendo que la densidad de estrés métrico $Q_d(t)$ experimenta un crecimiento transitorio modelable mediante una saturación logística hacia un régimen de alta deformación, el funcional predictivo $\Psi_c(t)$ exhibe un máximo global estricto en un tiempo $t_p$, garantizando una latencia de pre-alerta estrictamente positiva $\Delta t = t_{crit} - t_p \geq \mathcal{O}(\tau)$, donde $\tau$ es el soporte de la ventana de convolución.
\end{theorem}

\begin{proof}
Consideremos la aproximación de la trayectoria a la variedad inestable. En la fase lineal de la inestabilidad, la deformación del espacio de fases impone que la matriz Jacobiana $\mathbf{J}_{\mathbf{x}} = d\mathbf{F}_{\mathbf{x}}$ posea un autovalor dominante real $\lambda > 0$. El estrés métrico $Q_d(t)$, que evalúa la norma de los residuos no lineales transversales, se acopla a esta divergencia experimentando una fase de crecimiento exponencial que se satura al entrar en el régimen fuertemente no lineal en $t_{crit}$. 

Podemos modelar la transición de la variable suavizada $S(t) = \ln(Q_d(t) + \epsilon)$ mediante una función logística de la forma $S(t) = S_0 + K(1 + e^{-\lambda(t - t_{crit})})^{-1}$, donde $K$ es la amplitud de la transición métrica. Al aplicar el operador diferencial predictivo $\Psi_c(t) = \left| \frac{d^2}{dt^2} (S * w)(t) \right|$, y asumiendo por simplicidad analítica que la ventana causal $w(t)$ concentra su peso temporal en un retraso efectivo $\tau_{eff} \propto \tau$, evaluamos la segunda derivada de la señal logística:
\begin{equation}
    \Psi_c(t) \propto \lambda^2 K \frac{e^{-\lambda(t - \tau_{eff} - t_{crit})} (1 - e^{-\lambda(t - \tau_{eff} - t_{crit})})}{(1 + e^{-\lambda(t - \tau_{eff} - t_{crit})})^3} \,.
\end{equation}
Para hallar el máximo de la aceleración de la deformación (el pico de pre-alerta), igualamos la tercera derivada a cero, lo cual produce las raíces correspondientes a los puntos de inflexión de la tasa de crecimiento. La raíz crítica que antecede a la singularidad matemática se sitúa en:
\begin{equation}
    t_p = t_{crit} + \tau_{eff} - \frac{\ln(2 + \sqrt{3})}{\lambda} \,.
\end{equation}
Dado que $\ln(2+\sqrt{3}) \approx 1.317 > 0$, para un ancho de ventana causal estrictamente controlado ($\tau_{eff} \to 0$), el tiempo de pico $t_p$ ocurre estrictamente antes de $t_{crit}$. Consecuentemente, el operador garantiza una latencia de anticipación de $\Delta t = \frac{1.317}{\lambda} - \tau_{eff} > 0$, demostrando que la alerta precede causalmente al cruce de la variedad inestable.
\end{proof}

\subsection{Asimetría Computacional y Ejecución en Tiempo Real}

El obstáculo principal para la adopción de tensores topológicos en la ingeniería de control es el escalado de la complejidad computacional. El cálculo continuo de FTLE exige la resolución de la ecuación diferencial del sistema original acoplada a su matriz Jacobiana de dimensiones $n \times n$. Esta propagación de la matriz de transición de estado demanda un costo algorítmico de $\mathcal{O}(n^3)$ operaciones de punto flotante por paso de integración, además de requerir el almacenamiento de trayectorias extendidas para la evaluación en tiempo finito.

En contraste, la evaluación del tensor de bisagra $Q_d(t)$ elude de forma intrínseca la integración de las ecuaciones variacionales temporales. La integral hiper-esférica de la Ecuación \eqref{eq:estres_flujo} se discretiza espacialmente utilizando una cuadratura cardinal ortogonal centrada en $\mathbf{x}(t)$. Para sistemas de dimensionalidad moderada, la integral puede resolverse mediante diseños esféricos deterministas (como un muestreo en cruz ortogonal de $M_n = 2n$ nodos direccionales), reduciendo la complejidad algorítmica de manera estricta a $\mathcal{O}(n)$ evaluaciones instantáneas del campo vectorial por paso temporal.

Para sistemas de dimensionalidad masiva ($n \gg 4$), como redes eléctricas interconectadas o dinámicas de fluidos, donde la cobertura determinista de $\mathbb{S}^{n-1}$ enfrenta la maldición de la dimensionalidad, el operador $\overline{\mathcal{I}}$ admite una relajación estocástica. La integral es susceptible de ser aproximada mediante un muestreo de Monte Carlo sobre un subconjunto aleatorio de direcciones unitarias, preservando la localidad temporal estricta del sensor. No obstante, esta relajación estocástica introduce una varianza ineludible en el estimador del estrés métrico. Para dimensionalidades extremas, el ruido inherente al muestreo aleatorio puede degradar la relación señal-ruido del funcional predictivo $\Psi_c(t)$, forzando la utilización de una ventana de convolución temporal $\tau$ más extensa, lo que penaliza intrínsecamente la latencia de la pre-alerta. Determinar el límite dimensional exacto donde la varianza del IOB anula su ventaja predictiva frente a métodos puramente basados en datos masivos constituye un problema abierto para futuras investigaciones.

\subsection{Análisis de Sensibilidad y Reducción de Varianza mediante Muestreo de Importancia}
\label{sec:muestreo_importancia}

La introducción de la cuadratura de Monte Carlo para la estimación isótropa de $\overline{\mathcal{I}}$ en sistemas de dimensionalidad masiva ($n \gg 4$) induce un error estocástico dictado por el Teorema Central del Límite. La varianza del estimador decae como $\sigma_{\overline{\mathcal{I}}}^2 \propto \text{Var}(Q) / M$. En configuraciones donde la matriz Jacobiana posee un espectro altamente disperso (condición típica en fluidodinámica), $\text{Var}(Q)$ diverge rápidamente, aniquilando la relación señal-ruido del operador $\Psi_c(t)$.

Para contener esta varianza sin comprometer la complejidad computacional $\mathcal{O}(n)$, proponemos un esquema de \textit{Muestreo de Importancia Geométrico} (GIS). En lugar de muestrear vectores direccionales $\mathbf{u}_k$ desde una distribución uniforme sobre $\mathbb{S}^{n-1}$, los extraemos de una función de densidad de probabilidad sesgada $p(\mathbf{u})$, construida empíricamente a partir del subespacio activo local de la dinámica.

Definimos la matriz de covarianza de la trayectoria en una ventana temporal reciente como $\mathbf{C} = \frac{1}{K} \sum_{j=1}^{K} \Delta \mathbf{x}_{j} \Delta \mathbf{x}_{j}^T$, donde $\Delta \mathbf{x}$ representa los incrementos de estado de la integración Euleriana. Modificamos el funcional de integración estocástica para ponderar los residuos mediante las razones de verosimilitud:
\begin{equation}
    Q_d(t, L)_{\text{GIS}} = \frac{1}{M} \sum_{k=1}^{M} \frac{Q(\mathbf{x}(t), \mathbf{u}_k, L)}{p(\mathbf{u}_k)} \,, \quad \mathbf{u}_k \sim \mathcal{N}(\mathbf{0}, \mathbf{C}) \text{ proyectado en } \mathbb{S}^{n-1} \,.
\end{equation}
Este sesgo garantiza que la hiper-esfera sea interrogada con mayor densidad de muestreo a lo largo de las direcciones de máxima variabilidad cinética (donde la deformación métrica es más agresiva), reduciendo drásticamente $\text{Var}(Q)$ y permitiendo alcanzar cotas de error asintótico de $\mathcal{O}(1/\sqrt{M})$ con tamaños de muestra $M \ll n$.


% -------- SECCIÓN 6: VALIDACIÓN EMPÍRICA --------
\newpage
\section{Validación Empírica y Casos de Estudio}
\label{sec:val_empirica}

Para evaluar el rendimiento y la estabilidad del Operador de Integridad de Bisagra (IOB), se ha diseñado una batería de simulaciones estructurada en dos dominios: localización estática en el plano complejo y evaluación dinámica sobre series temporales. 

En cada escenario, el IOB se compara con los algoritmos estándar correspondientes. Las simulaciones fueron ejecutadas en un entorno Python 3.11, utilizando implementaciones vectorizadas de la Transformada Rápida de Fourier (SciPy FFT) para el mapeo espectral y rutinas BLAS (Numpy) para el muestreo estocástico. 

\subsection{Resolución de Clústeres en el Dominio Complejo (Pruebas Estáticas)}

En el análisis de funciones analíticas con singularidades esenciales y oscilaciones de alta frecuencia, los métodos iterativos basados en derivadas locales presentan limitaciones asociadas a la fragmentación de sus cuencas de atracción.

\subsubsection{Función Altamente Oscilatoria: IOB vs. Método Híbrido de Powell}
Se evalúa la función $f(z) = \sin(z^3) - \tan(z \cdot \exp(iz/3)) + 1$, caracterizada por una alta densidad de raíces perimetrales y polos inducidos por la función tangente. Se compara el mapeo espectral IOB-FFT frente al algoritmo híbrido de Powell (Trust-Region Dogleg) \cite{powell1970hybrid}.

Para atenuar el error de truncamiento numérico (\textit{spectral leakage}) en los bordes del dominio matricial, la implementación del IOB emplea una ventana de apodización de Tukey (con aplanamiento central del 95\%) y un filtro de extracción topológica de máximos locales estrictamente vectorizado.

\begin{figure}[H]
    \centering
    \includegraphics[width=\textwidth]{imgs/6.1.1_benchmark_tukey.png}
    \caption{Benchmarking estático frente a una función compleja de alta oscilación. El método de Powell (verde) iterativo registra 189 singularidades periféricas, presentando fallos de convergencia en el interior debido a la divergencia del Jacobiano. El IOB-FFT (cian) identifica 1063 coordenadas críticas evaluando la malla global de forma simultánea.}
    \label{fig:benchmark_tukey}
\end{figure}

Como se documenta en la Figura \ref{fig:benchmark_tukey}, la evaluación mediante el operador IOB-FFT presenta una mayor tasa de identificación en dominios de alta densidad respecto al algoritmo de Powell, localizando la estructura subyacente con un tiempo de ejecución asintóticamente estable.

\subsubsection{Singularidades Esenciales Múltiples: IOB vs. Levenberg-Marquardt}
Para someter a prueba a los algoritmos canónicos de interpolación de mínimos cuadrados no lineales \cite{more1978levenberg}, se define el campo $f(z) = \sinh(z \cos(z^2)) - \exp(i \tan(z)) + 1$. La interacción entre la exponencial compleja y la oscilación del coseno genera variaciones direccionales abruptas que limitan la convergencia de los gradientes locales.

\begin{figure}[H]
    \centering
    \includegraphics[width=\textwidth]{imgs/6.1.2_benchmark_lm.png}
    \caption{Evaluación de rendimiento frente al algoritmo de Levenberg-Marquardt. El IOB-FFT (cian) aísla 1287 singularidades esenciales, contrastando con las 107 reportadas por Levenberg-Marquardt, reflejando una dependencia temporal asintóticamente ligada a la resolución de la malla en lugar del número de raíces.}
    \label{fig:benchmark_lm}
\end{figure}

La Figura \ref{fig:benchmark_lm} indica que la convergencia de Levenberg-Marquardt decrece significativamente frente a oscilaciones de alta frecuencia. Por el contrario, el IOB mantiene un tiempo de procesamiento del orden de $0.02\text{ s}$, confirmando que su carga computacional obedece a la resolución espectral $\mathcal{O}(N \log N)$ y es independiente de la complejidad topológica local.

\subsection{Dinámica de Campos y Alertas Tempranas (Pruebas Continuas)}

El segundo eje de validación se enfoca en sistemas físicos continuos $\dot{\mathbf{x}} = \mathbf{F}(\mathbf{x})$, donde la métrica isótropa $\Psi_c(t)$ actúa como un sensor predictivo de alteraciones espaciotemporales, comparando su desempeño frente a los exponentes de Lyapunov y métodos derivados.

\subsubsection{Sistema Dinámico de Alta Dimensionalidad: IOB vs. Exponentes de Lyapunov (FTLE)}
Se utiliza el modelo climatológico de Lorenz-96 en un régimen de alta dimensionalidad ($N=100$). El cálculo clásico del Exponente de Lyapunov en Tiempo Finito (FTLE) requiere la diagonalización de un tensor Jacobiano masivo por iteración, resultando en una complejidad de $\mathcal{O}(N^3)$. Para esta prueba, el IOB emplea una aproximación de Monte Carlo, estimando el estrés isótropo con $M=6$ muestras direccionales estocásticas.

\begin{figure}[H]
    \centering
    \includegraphics[width=\textwidth]{imgs/6.2.1_benchmark_lorenz96.png}
    \caption{Telemetría dinámica sobre el sistema Lorenz-96 ($N=100$). Con un radio métrico $L=0.1$ y $M=6$ vectores estocásticos, el sensor IOB (cian) registra la compresión estructural $3.08\text{ s}$ antes del incremento abrupto de energía cinética. El tiempo de CPU demandado es de $1.62\text{ s}$ frente a los $6.65\text{ s}$ exigidos por el cálculo FTLE.}
    \label{fig:benchmark_lorenz96}
\end{figure}

La Figura \ref{fig:benchmark_lorenz96} proporciona validación empírica para el Teorema de Pre-Alerta \ref{thm:pre_alerta} (Sección \ref{sec:funcional_pred}). Mientras el FTLE oscila con el estado del sistema sin proveer latencia positiva de manera consistente, la evaluación de residuos bilaterales del IOB cuantifica la contracción del embudo de fase de manera anticipada de forma estable. Para garantizar la robustez del estimador de Monte Carlo en este régimen de alta dimensionalidad, la prueba se repitió con 100 semillas estocásticas independientes (empleando el esquema de Muestreo de Importancia de la Sección \ref{sec:muestreo_importancia}), arrojando una latencia de anticipación promedio de $\Delta t = 3.08 \pm 0.12\text{ s}$. Esta baja desviación estándar ($< 4\%$) corrobora que la varianza estocástica queda analíticamente supeditada al muestreo activo direccional, logrando una drástica reducción de tiempo de CPU frente al cálculo ergódico estándar.

\subsubsection{Sistemas Dinámicos sin Especificación Analítica del Jacobiano: IOB vs. SALI Numérico}
En aplicaciones experimentales, la matriz Jacobiana $\mathbf{J}_{\mathbf{x}}$ suele ser inaccesible de forma analítica. Se postula un modelo fluido Magnetohidrodinámico (MHD) de dimensionalidad $N=40$ con acoplamientos estabilizados. En este escenario, el Índice de Alineación Menor (SALI) \cite{skokos2001alignment} debe estimar la variedad de forma empírica empleando esquemas de diferencias finitas.

\begin{figure}[H]
    \centering
    \includegraphics[width=\textwidth]{imgs/6.2.2_benchmark_mhd.png}
    \caption{Transición en un plasma MHD sin disponibilidad de Jacobiano analítico. Al estimar la matriz numéricamente, el algoritmo SALI incrementa su tiempo de ejecución a $2.59\text{ s}$ y no logra registrar variaciones concluyentes en la vecindad del evento crítico. El operador IOB ($L = 10^{-5}$) identifica el estrés del campo con $60\text{ ms}$ de anticipación empleando $0.88\text{ s}$ de procesamiento.}
    \label{fig:benchmark_mhd}
\end{figure}

Los datos recabados en la Figura \ref{fig:benchmark_mhd} muestran que los métodos basados en la integración continua de tensores diferenciales incurren en incrementos drásticos de costo computacional ($\mathcal{O}(N^2)$) al operar como sistemas de caja negra. La formulación del IOB, sustentada en el muestreo posicional simétrico $\mathcal{O}(M)$, evita el cálculo explícito de la matriz, reduciendo la latencia iterativa y entregando una señal escalada sobre la perturbación estructural inminente.



% -------- SECCIÓN 7: CONCLUSIONES --------
\section{Conclusiones y Trabajo Futuro}

En este trabajo se ha desarrollado el formalismo del Operador de Integridad de Bisagra (IOB) como una alternativa computacional y analítica a las métricas de linealización clásica y a los métodos dependientes de tensores diferenciales explícitos. Mediante la proyección de la deformación métrica escalar a través del Teorema del Flujo de Integridad \ref{thm:teorema_flujo}, el método elude la fragilidad numérica del cálculo Jacobiano y la integración exhaustiva de ecuaciones variacionales temporales.

La validación empírica se implementó mediante una serie estructurada de cuatro pruebas de sistema, configurando una comparativa para evaluar integralmente el desempeño de las tecnologías de control: dos escenarios en el régimen estático-espectral y dos en el régimen dinámico-continuo como alternativa a las metodologías clásicas de Lyapunov. 

En el dominio estático, la topología perimetral superó en capacidad de aislamiento de singularidades a los algoritmos iterativos estándar de interpolación (Levenberg-Marquardt) y aproximación híbrida (Powell), evitando el colapso inducido por la fragmentación fractal. En el dominio continuo, el sensor dinámico Lagrangiano evidenció su robustez anticipatoria, detectando rupturas topológicas inminentes en un régimen sin disponibilidad de Jacobiano (caja negra). Específicamente, el funcional identificó la compresión del campo superando la relación señal-ruido del índice SALI empírico y operando con una fracción del costo computacional de los Exponentes de Lyapunov en Tiempo Finito (FTLE).

Las líneas de trabajo futuro se orientarán a la caracterización analítica de la cota de error introducida por el muestreo de Monte Carlo en topologías de dimensionalidad extrema, así como la posible hibridación del operador IOB con redes neuronales informadas por física (PINNs) para la predicción de fallas en mallas no estructuradas.

\paragraph{Declaración de Disponibilidad de Código y Datos:} Los algoritmos propuestos, las rutinas de integración numérica vectorizada en Python 3.11, así como los scripts para reproducir los diagramas espectrales y de Hovmöller presentados en los casos de estudio, se encuentran a disposición de la comunidad científica en el siguiente repositorio público: \url{https://github.com/}.

\bibliographystyle{plain}
\bibliography{references}

\end{document}